\subsection{Quiz}

To evaluate understanding of the Chain of Responsibility pattern, consider the following technical scenarios:

\subsubsection*{Question 1: Architectural Intent}
\textbf{Scenario:} You are designing a system where an incoming request must be processed by \textit{all} validation steps in sequence, regardless of whether a previous step succeeded or failed. Is standard Chain of Responsibility the optimal choice?

\begin{enumerate}[label=(\Alph*)]
    \item Yes, CoR is strictly designed to execute all handlers unconditionally.
    \item No, because standard CoR allows any handler to break the chain execution prematurely upon failure or completion.
    \item Yes, but only if implemented with multiple threads.
    \item No, Strategy pattern must be used instead.
\end{enumerate}

\textbf{Answer:} \textbf{(B)}. Standard CoR gives each handler the ability to stop propagation. If all steps must always execute, a pipeline or sequential processing model is more suitable.

\subsubsection*{Question 2: Runtime Behavior}
\textbf{Scenario:} Consider a chain composed of three handlers: \texttt{HandlerA} $\rightarrow$ \texttt{HandlerB} $\rightarrow$ \texttt{HandlerC}. If \texttt{HandlerB} processes a request and does not call the next handler, what is the outcome?

\begin{enumerate}[label=(\Alph*)]
    \item \texttt{HandlerC} will still execute automatically.
    \item A runtime exception is thrown.
    \item \texttt{HandlerC} is never invoked, and processing stops at \texttt{HandlerB}.
    \item The request is reprocessed from \texttt{HandlerA}.
\end{enumerate}

\textbf{Answer:} \textbf{(C)}. In CoR, continuation depends explicitly on invoking the next handler. If omitted, the chain terminates at that point.

\subsubsection*{Question 3: Advanced Design Consideration}
\textbf{Scenario:} You are implementing a logging system where multiple handlers must process the same request independently (e.g., writing to a file, sending alerts, and storing metrics), and no handler should block others from executing. Which design adjustment is most appropriate?

\begin{enumerate}[label=(\Alph*)]
    \item Use standard CoR without modification.
    \item Modify CoR so that each handler always forwards the request after processing.
    \item Replace CoR with an Observer (Publish-Subscribe) pattern.
    \item Use recursion to re-trigger the chain for each handler.
\end{enumerate}

\textbf{Answer:} \textbf{(C)}. When multiple independent consumers must react to the same event without affecting each other, the Observer pattern is a better fit. CoR is inherently sequential and control-flow driven, whereas Observer enables parallel, decoupled notifications.